Misinformation spreads inside groups of people who are both connected and like-minded,
and crosses into other groups through a small number of accounts and links. Most
containment methods choose where to intervene from network structure alone. This
project asks whether a second signal, the content a user posts, can identify
intervention targets that structure misses. We define, for every user, a
\emph{structural} community membership (from the social graph) and a \emph{content}
membership (what a classifier that reads only the user's text predicts that community
to be), and score the disagreement between the two. The classifier receives no edges
or graph-derived features as input; the graph communities are used only as its
training targets. The idea comes from
the use of complementary clusterings to locate ambiguous regions in active
learning~\citep{sani2026aas}. This mid-term report describes the motivation, related
work, the methodology as currently built, and the experimental setup on two real
social networks (TIMME political Twitter and the Pokec social network) and on
controlled synthetic networks, together with the current status and the plan for the
rest of the project. Results are not reported here.
